% Part:sets-functions-relations
% Chapter: size-of-sets
% Section: pairing-alt

\documentclass[../../../include/open-logic-section]{subfiles}

\begin{document}

\olfileid[ar]{sfr}{siz}{pai-alt}

\olsection{دالة مزاوجة بديلة}

\begin{explain}
ولـ$\Nat^{\OLMathNumeral{2}}$ تعدادات أخرى أيسر في استخراج المعكوس، منها هذا.
خذ قائمة الصحيحات الموجبة، واجعل بإزاء كل عدد مكانًا خاليًا.
ثم املأ الأمكنة تباعًا بالأزواج $\tuple{{\OLId{latin-ordinary}{n}},{\OLId{latin-ordinary}{m}}}$ بهذه الصنعة:
ابدأ بما أوله~${\OLMathNumeral{0}}$، أي $\tuple{{\OLMathNumeral{0}},{\OLId{latin-ordinary}{m}}}$؛ وضع أولها $\tuple{{\OLMathNumeral{0}},{\OLMathNumeral{0}}}$
في أول خال، واترك التالي، ثم ضع الثاني $\tuple{{\OLMathNumeral{0}},{\OLMathNumeral{1}}}$ في الخالي
الذي بعده، واترك آخر، وهكذا. فتحصل بداية التعداد غير المكتملة:
\[\small
\begin{array}{@{}c c c c c c c c c c c@{}}
\mathbf 1 & \mathbf 2 & \mathbf 3 & \mathbf 4 & \mathbf 5 & \mathbf 6 & \mathbf 7 & \mathbf 8 & \mathbf 9 & \mathbf{10} & \dots \\ \\
\tuple{{\OLMathNumeral{0}},{\OLMathNumeral{0}}} &  & \tuple{{\OLMathNumeral{0}},{\OLMathNumeral{1}}} &  & \tuple{{\OLMathNumeral{0}},{\OLMathNumeral{2}}} &  & \tuple{{\OLMathNumeral{0}},{\OLMathNumeral{3}}} & & \tuple{{\OLMathNumeral{0}},{\OLMathNumeral{4}}} &  & \dots \\
\end{array}
\]
ثم افعل بالأزواج $\tuple{{\OLMathNumeral{1}},{\OLId{latin-ordinary}{m}}}$ مثل ذلك في الأمكنة الباقية خالية،
فتملأ واحدًا وتدع واحدًا:
\[\small
\begin{array}{@{}c c c c c c c c c c c@{}}
\mathbf 1 & \mathbf 2 & \mathbf 3 & \mathbf 4 & \mathbf 5 & \mathbf 6 & \mathbf 7 & \mathbf 8 & \mathbf 9 & \mathbf{10} & \dots \\ \\
\tuple{{\OLMathNumeral{0}},{\OLMathNumeral{0}}} & \tuple{{\OLMathNumeral{1}},{\OLMathNumeral{0}}} & \tuple{{\OLMathNumeral{0}},{\OLMathNumeral{1}}} &  & \tuple{{\OLMathNumeral{0}},{\OLMathNumeral{2}}} & \tuple{{\OLMathNumeral{1}},{\OLMathNumeral{1}}} &
\tuple{{\OLMathNumeral{0}},{\OLMathNumeral{3}}} & & \tuple{{\OLMathNumeral{0}},{\OLMathNumeral{4}}} &  \tuple{{\OLMathNumeral{1}},{\OLMathNumeral{2}}} & \dots \\
\end{array}
\]
ثم أدخل $\tuple{{\OLMathNumeral{2}},{\OLId{latin-ordinary}{m}}}$ و$\tuple{{\OLMathNumeral{3}},{\OLId{latin-ordinary}{m}}}$ وما بعدها على القياس نفسه؛
فيصير أول التعداد المكتمل:
\[\small
\begin{array}{@{}cc c c c c c c c c c@{}}
\mathbf 1 & \mathbf 2 & \mathbf 3 & \mathbf 4 & \mathbf 5 & \mathbf 6 & \mathbf 7 & \mathbf 8 & \mathbf 9 & \mathbf{10} & \dots \\ \\
\tuple{{\OLMathNumeral{0}},{\OLMathNumeral{0}}} & \tuple{{\OLMathNumeral{1}},{\OLMathNumeral{0}}} & \tuple{{\OLMathNumeral{0}},{\OLMathNumeral{1}}} & \tuple{{\OLMathNumeral{2}},{\OLMathNumeral{0}}}  & \tuple{{\OLMathNumeral{0}},{\OLMathNumeral{2}}} &
\tuple{{\OLMathNumeral{1}},{\OLMathNumeral{1}}} & \tuple{{\OLMathNumeral{0}},{\OLMathNumeral{3}}} & \tuple{{\OLMathNumeral{3}},{\OLMathNumeral{0}}}  & \tuple{{\OLMathNumeral{0}},{\OLMathNumeral{4}}} &  \tuple{{\OLMathNumeral{1}},{\OLMathNumeral{2}}} & \dots \\
\end{array}
\]
\begin{editorial}
\textbf{تنبيه تصحيحي على الأصل (OLSIZ-003).}
سمّى الأصل الزوج الثاني $\tuple{{\OLMathNumeral{0}},{\OLMathNumeral{2}}}$، فأثبتناه $\tuple{{\OLMathNumeral{0}},{\OLMathNumeral{1}}}$؛
فهو الذي يقع في الموضع الثالث من جدول بداية التعداد الأول.
وكرر الأصل أسرة الأزواج $\tuple{{\OLMathNumeral{2}},{\OLId{latin-ordinary}{m}}}$، وصواب الثانية
$\tuple{{\OLMathNumeral{3}},{\OLId{latin-ordinary}{m}}}$؛ ويشهد له الزوج $\tuple{{\OLMathNumeral{3}},{\OLMathNumeral{0}}}$ في الموضع الثامن
من بداية التعداد المكتمل أعلاه.
\end{editorial}
وترقيم خانات المصفوفة بهذا التعداد لا يسلك الخط المتعرج السابق،
بل يعطي هذا الترتيب:
\[
\begin{array}{ c | c | c | c | c | c | c | c }
& \mathbf 0 & \mathbf 1 & \mathbf 2 & \mathbf 3 & \mathbf 4 & \mathbf 5 & \dots \\
\hline
\mathbf 0 & {\OLMathNumeral{1}} & {\OLMathNumeral{3}} & {\OLMathNumeral{5}} & {\OLMathNumeral{7}} & {\OLMathNumeral{9}} & {\OLMathNumeral{11}} & \dots \\
\hline
\mathbf 1 & {\OLMathNumeral{2}} & {\OLMathNumeral{6}} & {\OLMathNumeral{10}} & {\OLMathNumeral{14}} & {\OLMathNumeral{18}} & \dots & \dots \\
\hline
\mathbf 2 & {\OLMathNumeral{4}} & {\OLMathNumeral{12}} & {\OLMathNumeral{20}} & {\OLMathNumeral{28}} & \dots & \dots & \dots \\
\hline
\mathbf 3 & {\OLMathNumeral{8}} & {\OLMathNumeral{24}} & {\OLMathNumeral{40}} & \dots & \dots & \dots & \dots \\
\hline
\mathbf 4 & {\OLMathNumeral{16}} & {\OLMathNumeral{48}} & \dots & \dots & \dots & \dots & \dots \\
\hline
\mathbf 5 & {\OLMathNumeral{32}} & \dots & \dots & \dots & \dots & \dots & \dots \\
\hline
\vdots & \vdots & \vdots & \vdots & \vdots & \vdots & \vdots & \ddots\\
\end{array}
\]
فأزواج الصف~${\OLMathNumeral{0}}$ مواضعها فردية: موضع $\tuple{{\OLMathNumeral{0}},{\OLId{latin-ordinary}{m}}}$ هو ${\OLMathNumeral{2}}{\OLId{latin-ordinary}{m}}+{\OLMathNumeral{1}}$.
وأزواج الصف الثاني $\tuple{{\OLMathNumeral{1}},{\OLId{latin-ordinary}{m}}}$ مواضعها ضعف الفردي، أي
${\OLMathNumeral{2}} \cdot ({\OLMathNumeral{2}}{\OLId{latin-ordinary}{m}}+{\OLMathNumeral{1}})$. وأزواج الثالث $\tuple{{\OLMathNumeral{2}},{\OLId{latin-ordinary}{m}}}$ مواضعها أربعة
أمثاله، أي ${\OLMathNumeral{4}} \cdot ({\OLMathNumeral{2}}{\OLId{latin-ordinary}{m}}+{\OLMathNumeral{1}})$؛ ثم على هذا القياس. وعوامل
$({\OLMathNumeral{2}}{\OLId{latin-ordinary}{m}}+{\OLMathNumeral{1}})$ في الصفوف، أعني ${\OLMathNumeral{1}}$، ${\OLMathNumeral{2}}$، ${\OLMathNumeral{4}}$، ${\OLMathNumeral{8}}$، \dots، قوى~${\OLMathNumeral{2}}$:
${\OLMathNumeral{1}}= {\OLMathNumeral{2}}^{\OLMathNumeral{0}}$، ${\OLMathNumeral{2}} = {\OLMathNumeral{2}}^{\OLMathNumeral{1}}$، ${\OLMathNumeral{4}} = {\OLMathNumeral{2}}^{\OLMathNumeral{2}}$، ${\OLMathNumeral{8}} = {\OLMathNumeral{2}}^{\OLMathNumeral{3}}$، \dots\@
والأس هو المركبة الأولى للزوج؛ فعامل $\tuple{{\OLId{latin-ordinary}{n}},{\OLId{latin-ordinary}{m}}}$ هو ${\OLMathNumeral{2}}^{\OLId{latin-ordinary}{n}}$.
فتكون الصيغة الجامعة ${\OLMathNumeral{2}}^{\OLId{latin-ordinary}{n}} \cdot ({\OLMathNumeral{2}}{\OLId{latin-ordinary}{m}}+{\OLMathNumeral{1}})$. غير أن هذه القاعدة
تعطي صحيحات \emph{موجبة}؛ فـ$\tuple{{\OLMathNumeral{0}},{\OLMathNumeral{0}}}$ رتبته~${\OLMathNumeral{1}}$.
ولكي نبدأ بالرتبة~${\OLMathNumeral{0}}$ ننقص~${\OLMathNumeral{1}}$ من الحاصل، فنحصل على الآتي:
\end{explain}

\begin{ex}
خذ الدالة ${\OLId{latin-ordinary}{h}}\colon \Nat^{\OLMathNumeral{2}} \to \Nat$ ذات القاعدة
\[
{\OLId{latin-ordinary}{h}}({\OLId{latin-ordinary}{n}},{\OLId{latin-ordinary}{m}}) = {\OLMathNumeral{2}}^{\OLId{latin-ordinary}{n}} ({\OLMathNumeral{2}}{\OLId{latin-ordinary}{m}}+{\OLMathNumeral{1}}) - {\OLMathNumeral{1}}
\]
فهذه دالة مزاوجة لأزواج الطبيعيّات~$\Nat^{\OLMathNumeral{2}}$.
\end{ex}

\begin{explain}
وعلى هذا، في التعداد الثاني لـ$\Nat^{\OLMathNumeral{2}}$ يكون رمز الزوج
$\tuple{{\OLMathNumeral{0}},{\OLMathNumeral{0}}}$ هو ${\OLId{latin-ordinary}{h}}({\OLMathNumeral{0}},{\OLMathNumeral{0}}) = {\OLMathNumeral{2}}^{\OLMathNumeral{0}}({\OLMathNumeral{2}}\cdot {\OLMathNumeral{0}}+{\OLMathNumeral{1}}) - {\OLMathNumeral{1}} = {\OLMathNumeral{0}}$؛
ورمز $\tuple{{\OLMathNumeral{1}},{\OLMathNumeral{2}}}$ هو ${\OLMathNumeral{2}}^{{\OLMathNumeral{1}}} \cdot ({\OLMathNumeral{2}} \cdot {\OLMathNumeral{2}} + {\OLMathNumeral{1}}) - {\OLMathNumeral{1}} = {\OLMathNumeral{2}}
\cdot {\OLMathNumeral{5}} - {\OLMathNumeral{1}} = {\OLMathNumeral{9}}$؛ ورمز $\tuple{{\OLMathNumeral{2}},{\OLMathNumeral{6}}}$ هو ${\OLMathNumeral{2}}^{{\OLMathNumeral{2}}} \cdot ({\OLMathNumeral{2}}
\cdot {\OLMathNumeral{6}} + {\OLMathNumeral{1}}) - {\OLMathNumeral{1}} = {\OLMathNumeral{51}}$.
\end{explain}

وقد يكفي ترميز أزواج الطبيعيّات~$\Nat^{\OLMathNumeral{2}}$، ولا يشترط الشمول.
فمعكوسات هذه الترميزات تكون دوال جزئية فحسب.

\begin{ex}
خذ الدالة ${\OLId{latin-ordinary}{j}}\colon \Nat^{\OLMathNumeral{2}} \to \Nat^+$ ذات القاعدة
\[
{\OLId{latin-ordinary}{j}}({\OLId{latin-ordinary}{n}},{\OLId{latin-ordinary}{m}}) = {\OLMathNumeral{2}}^{\OLId{latin-ordinary}{n}}{\OLMathNumeral{3}}^{\OLId{latin-ordinary}{m}}
\]
فهذه دالة !!a{injective} $\Nat^{\OLMathNumeral{2}} \to \Nat$.
\end{ex}

\end{document}
